% .NET Aspire section

% 1 – What is .NET Aspire?
\QuestionSlide[\CategoryBadge[AspireColor!20]{.NET Aspire}]<1>{What is .NET Aspire?}

\begin{frame}[fragile]
  \frametitle{\AnswerTitle[\CategoryBadge[AspireColor!20]{.NET Aspire}]{What is .NET Aspire?}}

  {\footnotesize
  \textbf{.NET Aspire} is an opinionated, cloud-native stack introduced in .NET 8. It provides \textbf{orchestration} (run all services locally with one command), \textbf{service discovery} (services find each other by name), \textbf{pre-built integrations} (Redis, SQL, RabbitMQ), and a \textbf{developer dashboard} with traces and logs.
  }
\end{frame}

% 2 – App Host
\QuestionSlide[\CategoryBadge[AspireColor!20]{.NET Aspire}]<2>{What is the .NET Aspire App Host?}

\begin{frame}[fragile]
  \frametitle{\AnswerTitle[\CategoryBadge[AspireColor!20]{.NET Aspire}]{What is the Aspire App Host?}}

  {\footnotesize
  The \textbf{App Host} is a C\# project that acts as the \textbf{composition root} for your distributed app. It declares which services and infrastructure to start, wires up references between them, and launches all processes (including Docker containers) with a single \texttt{dotnet run}.
  }

  \begin{minted}[fontsize=\scriptsize]{csharp}
var builder = DistributedApplication.CreateBuilder(args);

var db = builder.AddPostgres("postgres")
               .AddDatabase("catalogdb");

var api = builder.AddProject<Projects.CatalogApi>("catalog-api")
                 .WithReference(db)
                 .WaitFor(db);

builder.AddProject<Projects.WebFrontend>("frontend")
       .WithReference(api);

await builder.Build().RunAsync();
  \end{minted}
\end{frame}

% 3 – Service Discovery
\QuestionSlide[\CategoryBadge[AspireColor!20]{.NET Aspire}]<2>{What is Service Discovery in .NET Aspire?}

\begin{frame}[fragile]
  \frametitle{\AnswerTitle[\CategoryBadge[AspireColor!20]{.NET Aspire}]{What is Service Discovery in .NET Aspire?}}

  {\footnotesize
  Aspire injects connection strings and URLs into each service as environment variables at runtime. \texttt{WithReference()} wires a producer to a consumer. \texttt{AddServiceDiscovery()} on the \texttt{HttpClient} resolves logical names like \texttt{https://catalog-api} to the actual address, enabling zero-config communication across environments.
  }

  \begin{minted}[fontsize=\scriptsize]{csharp}
// In the consumer's Program.cs
builder.Services.AddHttpClient<CatalogClient>(
    client => client.BaseAddress = new Uri("https+http://catalog-api"))
    .AddServiceDiscovery();

// URL "https+http://catalog-api" resolves to the actual host/port
// injected by Aspire — no hard-coded addresses anywhere.
  \end{minted}
\end{frame}

% 4 – Dashboard
\QuestionSlide[\CategoryBadge[AspireColor!20]{.NET Aspire}]<1>{What does the .NET Aspire Dashboard show?}

\begin{frame}[fragile]
  \frametitle{\AnswerTitle[\CategoryBadge[AspireColor!20]{.NET Aspire}]{What does the Aspire Dashboard show?}}

  {\footnotesize
  The built-in \textbf{Aspire Dashboard} (launched automatically during local development) provides:
  \begin{itemize}
    \item \textbf{Resources} – all running projects and containers with health state.
    \item \textbf{Distributed Traces} – OpenTelemetry traces spanning multiple services.
    \item \textbf{Structured Logs} – aggregated logs from all services.
    \item \textbf{Metrics} – resource consumption and custom app metrics.
  \end{itemize}
  In production, forward the same telemetry to Azure Monitor, Grafana, or Jaeger.
  }
\end{frame}

% 5 – Aspire Integrations
\QuestionSlide[\CategoryBadge[AspireColor!20]{.NET Aspire}]<2>{What are .NET Aspire Integrations?}

\begin{frame}[fragile]
  \frametitle{\AnswerTitle[\CategoryBadge[AspireColor!20]{.NET Aspire}]{What are .NET Aspire Integrations?}}

  {\footnotesize
  \textbf{Integrations} are NuGet packages that wire up infrastructure components (Redis, PostgreSQL, Azure Service Bus, RabbitMQ, etc.) into the Aspire pipeline. \textbf{Hosting integrations} add resources to the App Host; \textbf{client integrations} configure the connection inside consuming services via DI — both with OpenTelemetry and health checks pre-wired.
  }

  \begin{minted}[fontsize=\scriptsize]{csharp}
// App Host – hosting integration
var cache = builder.AddRedis("cache");

// Consumer project – client integration
builder.AddRedisClient("cache");  // registers IConnectionMultiplexer
// or
builder.AddRedisDistributedCache("cache"); // registers IDistributedCache
  \end{minted}
\end{frame}

% 6 – WaitFor & health readiness
\QuestionSlide[\CategoryBadge[AspireColor!20]{.NET Aspire}]<2>{What does WaitFor do in .NET Aspire?}

\begin{frame}[fragile]
  \frametitle{\AnswerTitle[\CategoryBadge[AspireColor!20]{.NET Aspire}]{What does WaitFor do in .NET Aspire?}}

  {\footnotesize
  \texttt{.WaitFor(resource)} makes the App Host delay starting a dependent service until the resource reports \textbf{healthy}. This replaces fragile sleep-based startup scripts in \texttt{docker-compose}. For containers, health is determined by Docker health checks; for databases, by a successful connection probe.
  }

  \begin{minted}[fontsize=\scriptsize]{csharp}
var postgres = builder.AddPostgres("pg").AddDatabase("mydb");

// API won't start until "pg" container is healthy
builder.AddProject<Projects.MyApi>("my-api")
       .WithReference(postgres)
       .WaitFor(postgres);
  \end{minted}
\end{frame}

% 7 – Aspire Manifest & Deployment
\QuestionSlide[\CategoryBadge[AspireColor!20]{.NET Aspire}]<3>{What is the Aspire Manifest and how is it used for deployment?}

\begin{frame}[fragile]
  \frametitle{\AnswerTitle[\CategoryBadge[AspireColor!20]{.NET Aspire}]{What is the Aspire Manifest?}}

  {\footnotesize
  Running the App Host with \texttt{--publisher manifest} outputs a \textbf{JSON manifest} describing all resources and their relationships. Tools like \textbf{azd} (Azure Developer CLI) and \textbf{Aspir8} (Kubernetes) read this manifest to generate deployment-ready \texttt{bicep}, \texttt{Helm} charts, or Docker Compose files — keeping deployment config derived from code.
  }

  \begin{minted}[fontsize=\scriptsize]{bash}
# Generate manifest
dotnet run --project AppHost \
    --publisher manifest \
    --output-path aspire-manifest.json

# Deploy to Azure Container Apps via azd
azd up  # reads manifest and provisions infrastructure
  \end{minted}
\end{frame}

% 8 – Aspire vs docker-compose
\QuestionSlide[\CategoryBadge[AspireColor!20]{.NET Aspire}]<2>{How does .NET Aspire differ from docker-compose?}

\begin{frame}[fragile]
  \frametitle{\AnswerTitle[\CategoryBadge[AspireColor!20]{.NET Aspire}]{Aspire vs docker-compose?}}

  {\footnotesize
  \begin{tabular}{ll}
    \textbf{docker-compose} & \textbf{.NET Aspire} \\
    YAML config & C\# code (type-safe) \\
    Manual service URLs & Automatic service discovery \\
    No observability & Dashboard with traces/logs built in \\
    Separate health checks & Health checks auto-wired \\
    No IDE integration & F5 debug across all services \\
  \end{tabular}
  \vspace{0.5em}
  Aspire can \emph{also} launch Docker containers via \texttt{AddContainer()}, but adds the orchestration layer on top.
  }
\end{frame}
